\section{Release Construction and Validation}
\label{sec:release-validation}

The release packages existing generated candidate rows into a benchmark
artifact without requiring raw traces, and validates that package before
any number in this paper is computed from it. Validation checks schema
completeness, split normalization, presence of the selected families and
absence of blocked ones, duplicate candidates within a decision,
consistency of decision metadata, decision-view count consistency, and
checksums (Table~\ref{tab:artifact-layout} summarizes the package
components these checks cover). This validation is itself part of the
contribution: a released benchmark is only as trustworthy as the checks
that produced it, so we treat it as reproducible evidence rather than an
implementation detail.

\begin{sloppypar}
The evaluated open-release snapshot passed metadata repair, the full
verification suite, publication-bundle validation, and end-to-end release
validation, all via reproducible scripts rather than manual inspection.
This is validation of the release package itself; it does not claim that
public hosting or DOI assignment is complete, and superseded status notes
are never treated as evidence in their own right.
\end{sloppypar}

\begin{table}[t]
  \centering
  \footnotesize
  \caption{Release-package layout using release-relative paths for the evaluated benchmark snapshot.}
  \label{tab:artifact-layout}
  \begin{tabularx}{\columnwidth}{@{}P{0.31\columnwidth}L@{}}
    \toprule
    artifact & relative\_path \\
    \midrule
    candidate rows & data/candidate\_rows/ \\
    decision view & data/decision\_view/decision\_view.parquet \\
    pairwise sample & data/pairwise\_sample/pairwise\_sample.parquet \\
    release manifest & metadata/release\_manifest.json \\
    checksums & metadata/checksums.sha256 \\
    \bottomrule
  \end{tabularx}
\end{table}


\subsection{Provenance of the evaluation results in this paper}

Beyond the released dataset artifact above, every result reported in
Sections~\ref{sec:characterization}--\ref{sec:continuation} traces to an
independently frozen, gate-validated analysis artifact, each with its own
committed scripts, provenance record, and (where applicable) an
independent second code path that recomputes headline numbers and
requires exact agreement: the target-discriminativeness audit, the Tier-1
closed-loop production evidence (14/14 validity gates; a byte-verified
evidence freeze recording the SHA256 of every output file both before and
after freezing), the offline$\leftrightarrow$closed-loop linkage analysis,
the mechanistic workload analysis, and the continuation-policy sensitivity
pilot and full sampled study (prelaunch LRU-equivalence gate;
post-completion validity gates; an independent recheck script requiring
$10^{-9}$ numerical agreement with the primary analysis), plus the
full-population MRU continuation census (60/60 chunks, 2{,}363{,}286
decision-horizon records, 30/30 validity gates, independent recheck
passed). All six are scripted, reproducible, and checked into the
repository alongside this manuscript as compact evidence artifacts.

\textbf{Two label sources must not be conflated.} The canonical labels
evaluated throughout this paper always use LRU continuation
(Section~\ref{sec:label-generation}). The MRU and random \emph{relabelings}
computed for the continuation-sensitivity study
(Section~\ref{sec:continuation}) exist only to measure sensitivity and are
never treated elsewhere as the canonical label.

\textbf{The population-scale MRU continuation census is frozen evidence.}
The census (run \texttt{20260914T042528Z\_1a29e773a113}) is complete,
validated, and independently rechecked; its manifest and validation report
are committed to the repository. The raw 1.6 GB census output is preserved
separately outside this
manuscript branch; the manuscript imports only compact summaries,
validation reports, manifests, and Figure~\ref{fig:continuation-robustness}
source data.

\textbf{A known, unresolved provenance limitation.} The exact historical
commit provenance of the full-release label generator is incomplete. This
does not affect the validity of the frozen analysis artifacts above, each
of which independently
records its own generation provenance, but it means the full historical
lineage of the generator code itself predating those artifacts is not
completely reconstructable from repository history alone.
